\begin{table}
\centering
\caption{Input ablation using locked experiment outputs. Recent gauge history provides most of the one-hour correction skill, while retaining NWM adds 12.0--44.3 percentage points beyond USGS+ERA5-Land. RMSE is reported in m$^3$ s$^{-1}$.}
\small
\resizebox{\linewidth}{!}{%
\begin{tabular}{l cc cc cc}
\toprule
\textbf{Site} & \multicolumn{2}{c}{\textbf{Gauge-free surrogate}} & \multicolumn{2}{c}{\textbf{USGS+ERA5-Land}} & \multicolumn{2}{c}{\textbf{USGS+NWM+ERA5-Land}} \\
\cmidrule(lr){2-3}\cmidrule(lr){4-5}\cmidrule(lr){6-7}
 & RMSE & $\Delta$RMSE & RMSE & $\Delta$RMSE & RMSE & $\Delta$RMSE \\
\midrule
Jefferson, NC & 14.34 & 8.5\% & 10.45 & 33.4\% & \textbf{8.57} & \textbf{45.3\%} \\
Galax, VA & 60.61 & 21.8\% & 57.53 & 25.8\% & \textbf{40.15} & \textbf{48.2\%} \\
Sugar Grove, NC & 7.53 & 5.0\% & 7.00 & 11.7\% & \textbf{3.48} & \textbf{56.0\%} \\
\bottomrule
\end{tabular}
}
\label{tab:input_ablation}
\end{table}
